\section{模型的假设}
\begin{enumerate}
  \item 每个十分钟区间内功率近似恒定，控制器能够响应当段已实现负荷和光伏。题目未给出额外响应延迟、输电限制、老化费用或自放电，本文不另行编造这些参数。
  \item 充、放电效率分别取0.9，因而往返效率为0.81；储能的充放电功率在交流母线侧计量。将90\%解释为往返效率的情形另作敏感性比较。
  \item 不售电，外网购电没有题给的功率上限。允许不利用无法吸收的光伏或不接收已经付费的计划电量，其购电支出照付。两类余量统一称为未利用电量，不全部称为弃光。
  \item 问题一执行日循环等式。问题二至四从1月1日6000\,kWh连续运行；规划器使用日末预测储电量至少6000\,kWh的库存恢复设计，实际执行不强制达到这一目标，也不在午夜免费重置。
  \item 波动电价的未来值不视为提前公布。规划使用历史预测价，执行和结算使用对应交付段已实现的价格；主解释中的“交易时刻”取该交付区间。
\end{enumerate}
